\section{Diskussion}
\label{sec:Diskussion}

\subsection{Winkelrichtgröße $D$ und Eigenträgheitsmoment $I_D$}

Die Bestimmung der Winkelrichtgröße über die lineare Regression ergab $D = (0{,}0241 \pm 0{,}0009)\,\unit{\newton\meter\per\radian}$ bei einer akzeptablen relativen Unsicherheit von rund $3{,}7\,\%$. Im zugehörigen Plot (Abbildung~\ref{fig:d}) zeigt sich allerdings eine leichte systematische Abweichung der Messpunkte von der Ursprungsgeraden: Während die Kräfte bei kleinen Auslenkwinkeln leicht oberhalb der Ausgleichsgeraden liegen, flacht der Verlauf bei großen Winkeln ($\varphi \geq \qty{90}{\degree}$) ab. Dies lässt sich auf Probleme mit der Federwaage sowie der Vorrichtungsklemme der Drillachse zurückführen. Aufgrund von mehreren Problemen mit der neutralen Einstellung der Federwaage wurde diese oftmals gewechselt, jedoch flossen auch eventuelle innere Reibung oder mögliche Nichtlinearitäten der Feder in den Fehler mit ein. Dazu wird vermutet, dass die Vorrichtungsklemme nicht fest genug angezogen werden konnte, da es, wie in Tabelle~\ref{tab:d} zu sehen, zu zwei gleichen Werten für $F$ bei unterschiedlichen $\phi$ gekommen ist. Daher ist anzunehmen, dass auch weiterhin Verrutschungen vorgekommen sind.

Beim Nachweis des Satzes von Steiner (Abbildung~\ref{fig:steiner}) stimmt die experimentell ermittelte Geradensteigung $s = (813 \pm 20)\,\unit{\second\squared\per\meter\squared}$ innerhalb von rund $5\,\%$ (etwa $2\sigma$) gut mit dem theoretisch erwarteten Wert $s_{\text{theo}} = \frac{8\pi^2 m}{D} \approx \qty{855}{\second\squared\per\meter\squared}$ überein. Das daraus abgeleitete Eigenträgheitsmoment der Drillachse $I_D = (0{,}00072 \pm 0{,}00047)\,\unit{\kilogram\meter\squared}$ weist hingegen eine große relative Unsicherheit von etwa $65\,\%$ auf. Ursache hierfür ist die Differenzbildung $I_D = I_0 - I_s$, bei der das theoretische Stabträgheitsmoment $I_s = \qty{0,00289}{\kilogram\meter\squared}$ vom gesamten Restträgheitsmoment $I_0 = (0{,}00361 \pm 0{,}00047)\,\unit{\kilogram\meter\squared}$ abgezogen wird, während die volle absolute Unsicherheit aus dem Achsenabschnitt $t$ erhalten bleibt. Zudem enthält dieses $I_D$ sowohl das vernachlässigte Eigenträgheitsmoment der beiden verschobenen Zylinder um ihre eigene Schwerpunktsachse als auch das Trägheitsmoment der Klemmvorrichtung, mit der die Metallstange auf der Drillachse befestigt wurde.

\subsection{Trägheitsmomente der einfachen Körper und der unregelmäßigen Puppe}

Bei der Untersuchung der Einzelkörper liegen die experimentellen Werte für die Holzkugel ($I_{\text{K, 0}} = (0{,}00004 \pm 0{,}00047)\,\unit{\kilogram\meter\squared}$) und den vertikalen Metallzylinder ($I_{\text{Z, 0}} = (0{,}00022 \pm 0{,}00047)\,\unit{\kilogram\meter\squared}$) deutlich unterhalb der theoretischen Erwartungswerte ($I_{\text{K, theo}} = \qty{0,00097}{\kilogram\meter\squared}$ bzw. $I_{\text{Z, theo}} = \qty{0,00072}{\kilogram\meter\squared}$). Da die Messunsicherheit $u_{I_0}) = \qty{0,00047}{\kilogram\meter\squared}$ jeweils größer als der Messwert selbst ist (relative Unsicherheit $> 100\,\%$) und somit selbst die Null im $1\sigma$-Intervall liegt, besitzen die absoluten Einzelwerte nur eine geringe statistische Aussagekraft.

Noch deutlicher tritt dieser Effekt bei der Holzpuppe auf: Hier führt der Abzug von $I_D$ in beiden Haltungen zu unphysikalisch negativen Werten ($I_{\text{Pos1, 0}} = (-0{,}00025 \pm 0{,}00047)\,\unit{\kilogram\meter\squared}$ und $I_{\text{Pos2, 0}} = (-0{,}00045 \pm 0{,}00047)\,\unit{\kilogram\meter\squared}$), während das Zylindermodell theoretisch $I_{\text{Pos1, theo}} = (0{,}00058 \pm 0{,}00003)\,\unit{\kilogram\meter\squared}$ und $I_{\text{Pos2, theo}} = (0{,}00025 \pm 0{,}00002)\,\unit{\kilogram\meter\squared}$ vorhersagt. Der Hauptgrund für diese systematische Unterschätzung liegt im unterschiedlichen Versuchsaufbau: Während $I_D$ im Steiner-Versuch inklusive der schweren Metallstange bestimmt wurde, wurden Kugel, Zylinder und Puppe direkt ohne diese Halterung auf die Drillachse aufgesetzt. Da das Trägheitsmoment der Puppe kleiner als das der Klemmhalterung ist, entsteht durch die Subtraktion von $I_D$ ein negativer Offset.

Bildet man hingegen die Differenz $\Delta I = I_{\text{Pos1}} - I_{\text{Pos2}}$ zwischen beiden Positionen der Puppe, kürzt sich der systematische Fehler durch $I_D$ exakt heraus und die Unsicherheit sinkt um mehr als eine Größenordnung. Der experimentelle Wert $\Delta I_0 = (0{,}00020 \pm 0{,}00003)\,\unit{\kilogram\meter\squared}$ bestätigt qualitativ die theoretisch erwartete Zunahme von $\Delta I_{\text{theo}} = (0{,}00033 \pm 0{,}00002)\,\unit{\kilogram\meter\squared}$ beim Ausstrecken der Beine. Die verbleibende Abweichung zwischen Theorie und Experiment lässt sich auf mehrere Modell- und Messgrenzen zurückführen:
\begin{itemize}
	\item \textbf{Geometrische Vereinfachungen:} Die Annahme homogener Vollzylinder mit konstanter Dichte vernachlässigt die verjüngte Form der Gliedmaßen, Gelenkzwischenräume, das Bohrloch in der Holzkugel sowie mögliche Dichteinhomogenitäten im Holz.
	\item \textbf{Ausrichtung der Puppe:} Bereits geringe Abweichungen vom idealen $\qty{90}{\degree}$-Winkel beim Ausstrecken der Beine verringern den effektiven Schwerpunktsabstand $a_2$ quadratisch und reduzieren damit das reale Trägheitsmoment gegenüber der Theorie.
	\item \textbf{Zeitmessung:} Da ein Mensch im Mittel eine rein visuelle Reaktionszeit von 200 bis 300 Millisekunden besitzt, erhöht sich das Risiko, größere Fehler bei kürzeren Perioden zu machen. Bei den sehr kurzen Periodendauern von unter einer Sekunde ($T \approx \qty{0,60}{\second}$ bis $\qty{0,94}{\second}$) wirken sich menschliche Reaktionszeiten beim manuellen Stoppen einzelner Schwingungen prozentual deutlich stärker aus als im Steiner-Versuch.
\end{itemize}